\documentclass[11pt,reqno]{amsart}

\usepackage[T1]{fontenc}
\usepackage[utf8]{inputenc}
\usepackage{amsmath,amssymb,amsthm}
\usepackage{mathtools}
\usepackage[margin=1.15in]{geometry}
\usepackage{booktabs}
\usepackage{enumitem}
\usepackage[colorlinks=true,linkcolor=blue,citecolor=blue,urlcolor=blue]{hyperref}

%%% theorem environments
\theoremstyle{plain}
\newtheorem{theorem}{Theorem}
\newtheorem{lemma}[theorem]{Lemma}
\newtheorem{proposition}[theorem]{Proposition}
\newtheorem{corollary}[theorem]{Corollary}
\newtheorem{conjecture}[theorem]{Conjecture}
\newtheorem{fact}[theorem]{Fact}

\theoremstyle{definition}
\newtheorem{definition}[theorem]{Definition}
\newtheorem{remark}[theorem]{Remark}
\newtheorem{question}[theorem]{Question}

\numberwithin{theorem}{section}

%%% macros
\newcommand{\cube}{\{-1,1\}}
\newcommand{\bits}{\{-1,0,1\}}
\newcommand{\Inf}{\mathbf{Inf}}
\newcommand{\Ex}{\mathbf{E}}
\newcommand{\Prb}{\mathbf{Pr}}
\newcommand{\sd}{\mathbin{\triangle}}
\newcommand{\lk}{\operatorname{lk}}
\newcommand{\eps}{\varepsilon}
\newcommand{\bbF}{\mathbb{F}}
\newcommand{\bbZ}{\mathbb{Z}}
\newcommand{\bbR}{\mathbb{R}}
\newcommand{\nn}{n}
\newcommand{\Fhat}[1]{\widehat{#1}}
\newcommand{\XiC}{\Xi}
\newcommand{\todo}[1]{\textup{\textbf{[\,todo: #1\,]}}}
\DeclareMathOperator{\Deg}{deg}

\begin{document}

\title[The Nisan--Szegedy bound is never tight, and $R_3=10$]
{The Nisan--Szegedy bound on relevant variables\\ is never tight, and $R_3 = 10$}

\author{Andrew Moffat}
\thanks{The mathematics in this paper was developed with AI assistance (Claude).
All statements have been independently refereed inside the project and, where indicated,
verified by computer; responsibility for any remaining error is the author's.}

\date{\today}

\subjclass[2020]{06E30, 68Q06, 05E99}
\keywords{Boolean functions, Fourier degree, relevant variables, juntas, Nisan--Szegedy bound}

\begin{abstract}
Nisan and Szegedy proved in 1994 that a Boolean function whose real multilinear degree is at
most $d$ depends on at most $d\,2^{d-1}$ variables. Writing $R_d$ for the maximum number of
relevant variables of a degree-$d$ Boolean function, this says $R_d \le d\,2^{d-1}$, and the
bound is attained for $d = 1$ and $d = 2$. We prove two results.
First, the Nisan--Szegedy bound is never attained for $d \ge 3$: $R_d \le d\,2^{d-1} - 1$ for
every $d \ge 3$. The proof shows that in a hypothetical extremal function every derivative is
a character times the indicator of an affine subspace, that every vertex link is the facet
complex of a cross-polytope, and that the resulting term hypergraph splits into at least two
cross-polytope boundaries on disjoint variable sets, which no Boolean function admits.
Second, we determine $R_3$ exactly: $R_3 = 10$, attained by the Chiarelli--Hatami--Saks
function $\XiC_3$. Consequently $R_1 = 1$, $R_2 = 4$, $R_3 = 10$, and $22 \le R_4 \le 31$.
The value $R_3 = 10$ is confirmed independently by an exhaustive canonical-form search over
a finite reformulation of the problem, and the first structural steps of the argument have
been formalised in Lean~4.
\end{abstract}

\maketitle

\section{Introduction}\label{sec:intro}

\subsection{The problem}
Let $f \colon \cube^n \to \cube$ be a Boolean function, written in the $\pm 1$ convention, and
let
\[
  f = \sum_{S \subseteq [n]} \Fhat{f}(S)\, \chi_S,
  \qquad \chi_S(x) = \prod_{i \in S} x_i,
\]
be its Fourier (real multilinear) expansion. The \emph{degree} $\Deg f$ is the largest $|S|$
with $\Fhat{f}(S) \ne 0$. A variable $x_i$ is \emph{relevant} if $f$ depends on it, that is,
if $f(x) \ne f(x^{\oplus i})$ for some $x$, where $x^{\oplus i}$ flips the $i$-th coordinate;
equivalently, if $\Fhat{f}(S) \ne 0$ for some $S \ni i$.

Set
\[
  R_d := \sup \{\, n : \text{some } f \colon \cube^n \to \cube \text{ with } \Deg f \le d
  \text{ has all } n \text{ variables relevant} \,\}.
\]
That $R_d$ is finite is the theorem of Nisan and Szegedy \cite{NS94}: a degree-$d$ Boolean
function is a $d\,2^{d-1}$-junta, so
\begin{equation}\label{eq:ns}
  R_d \;\le\; d\,2^{d-1}.
\end{equation}
Chiarelli, Hatami and Saks \cite{CHS20} showed that the truth is exponential rather than
$d\,2^{d-1}$: there is an absolute constant $C \le 6.614$ with $R_d \le C\, 2^d$, and they
gave, for each $d$, an explicit degree-$d$ function $\XiC_d$ with
\begin{equation}\label{eq:chslower}
  3 \cdot 2^{d-1} - 2 \quad \text{relevant variables},
\end{equation}
so that $R_d \ge 3 \cdot 2^{d-1} - 2$. Wellens \cite{Wel22} improved the constant to $4.394$.
Writing
\[
  C_{44} := \sup_{d \ge 1} \frac{R_d}{2^d},
\]
the current state of the art is $3/2 \le C_{44} \le 4.394$; this is entry~44a on Tao's
optimisation problems page \cite{TaoOP44a}.

The improved constants of \cite{CHS20,Wel22} are asymptotic. For small $d$ they are worse than
\eqref{eq:ns}: the Chiarelli--Hatami--Saks optimisation attains its minimum at $d = 12$, and
Wellens' bound beats $d\,2^{d-1}$ only from roughly $d \ge 9$ on. For $d \le 8$, therefore,
\eqref{eq:ns} is still the best published upper bound, and the exact values of $R_d$ have not
been recorded anywhere for any $d \ge 3$: the problem page \cite{TaoOP44a} lists
$10 \le R_3 \le 12$ and states that exact values are not known.

\subsection{Results}
We settle the tightness of \eqref{eq:ns} for every $d$, and determine $R_3$.

\begin{theorem}\label{thm:main-ns}
For every $d \ge 3$, no Boolean function of degree $d$ has $d\,2^{d-1}$ relevant variables.
Hence
\[
  R_d \;\le\; d\,2^{d-1} - 1 \qquad (d \ge 3).
\]
\end{theorem}

\begin{theorem}\label{thm:main-r3}
$R_3 = 10$. That is, a Boolean function of degree $3$ has at most $10$ relevant variables, and
$10$ is attained by the Chiarelli--Hatami--Saks function $\XiC_3$.
\end{theorem}

Theorem~\ref{thm:main-ns} is sharp in the sense that $d \ge 3$ cannot be weakened: for $d = 1$
the function $f(x) = x_1$ has $1 = 1 \cdot 2^0$ relevant variable, and for $d = 2$ the function
\begin{equation}\label{eq:degree2extremal}
  f(x_1,x_2,x_3,x_4) = \tfrac12\left( x_1x_2 + x_1x_3 + x_2x_4 - x_3x_4 \right)
\end{equation}
is Boolean of degree $2$ with $4 = 2 \cdot 2^1$ relevant variables. So the Nisan--Szegedy bound
is attained exactly for $d \in \{1,2\}$.

Combining the two theorems with \eqref{eq:chslower} and \eqref{eq:ns}, the exact state of
knowledge for small $d$ is now
\[
  R_1 = 1, \qquad R_2 = 4, \qquad R_3 = 10, \qquad 22 \le R_4 \le 31,
\]
the upper bound $31$ for $d = 4$ being $4 \cdot 2^3 - 1$ from Theorem~\ref{thm:main-ns}.

\subsection{Method}
Both theorems run on the same structural device, which we call the \emph{link of a derivative}.
If $f$ has degree $d$ and $x_i$ is relevant, then the discrete derivative
\[
  D_i f(x) := \frac{f(x^{i \to 1}) - f(x^{i \to -1})}{2}
  = \sum_{S \ni i} \Fhat{f}(S)\, \chi_{S \setminus i}
\]
is a nonzero $\bits$-valued function of degree at most $d - 1$, and by granularity all of its
Fourier coefficients are integer multiples of $2^{1-d}$. Being $\bits$-valued is a very rigid
constraint on a sparse, granular Fourier expansion. Lemma~\ref{lem:A} makes this precise at the
extreme end: if the Fourier mass of $D_i f$ is as small as it can possibly be, then the support
of $D_i f$ is an affine subspace of $\bbF_2^n$ and all its coefficients are $\pm 2^{1-d}$, i.e.\
$D_i f$ is a character times the indicator of an affine subspace of codimension $d-1$. A
counting argument (Lemma~\ref{lem:B}) turns the affine subspace into the facet family of a
cross-polytope when all the support sets have the same size, and a closure argument propagates
the cross-polytope from one vertex to its neighbours until the whole term hypergraph is a
disjoint union of $2^{d-2}$ cross-polytope boundaries. For $d \ge 3$ there are at least two of
them, and a sum of two non-constant functions on disjoint variable sets is never $\pm1$-valued
(Lemma~\ref{lem:C}). That is Theorem~\ref{thm:main-ns}.

For Theorem~\ref{thm:main-r3} the same device is pushed through a finer case analysis. Part~I
($R_3 \le 11$, Section~\ref{sec:partI}) is the specialisation of the above to $d = 3$, where
the cross-polytope is an octahedron and the term hypergraph is a closed triangulated surface
with all vertex degrees $4$. Part~II ($R_3 \le 10$, Section~\ref{sec:partII} and
Appendix~\ref{app:cases}) has an extra unit of slack in the influence budget, which allows the
function to carry terms of degree below the top level and allows two vertices to have larger
derivative mass. The slack is accounted for by a single identity (Section~\ref{sec:bookkeep}),
and the resulting finitely many configurations are eliminated one by one.

\subsection{Verification}
The main theorems are proved by hand. Independently of the proofs, we reformulate
``no degree-$3$ Boolean function has $n$ relevant variables'' as a purely finite arithmetic
statement $F(n)$ about integer vectors (Section~\ref{sec:finite}), and run an exhaustive
orderly search over canonical forms. The search finds no solution for $n = 11$ and exactly one
orbit for $n = 10$, namely $\XiC_3$. A separate SAT-based structure search by an independent
referee finds the same short list of candidate term structures and no Boolean sign vector on
any of them. Finally, the first structural steps of Part~I have been formalised in Lean~4 and
kernel-checked on the standard axioms; we are explicit in Section~\ref{sec:lean} about what is
and what is not machine-checked (the finite statements $F(11)$ and $F(12)$ are \emph{not}).

\subsection{Related work and novelty}
Nisan and Szegedy \cite{NS94} do not discuss whether \eqref{eq:ns} is tight for $d \ge 3$.
Chiarelli--Hatami--Saks \cite{CHS20} and Wellens \cite{Wel19,Wel22} beat $d\,2^{d-1}$ only for
large $d$, and neither states an exact value of $R_3$ or $R_4$ nor addresses tightness at small
$d$; Wellens' Table~1 tabulates linear-programming bounds on block sensitivity, not on the
number of relevant variables. A literature sweep carried out for this project found no paper,
preprint, thesis or problem-tracker entry recording an exact value of $R_d$ for any $d \ge 3$.
The acknowledgements of \cite{CHS20} thank Avishay Tal for sharing exhaustive-search code for
Boolean functions on few variables, but no results of such a search are reported there. To the
best of our knowledge, Theorems~\ref{thm:main-ns} and \ref{thm:main-r3} are new.

\subsection{Organisation}
Section~\ref{sec:prelim} fixes notation and collects the standard facts we use.
Section~\ref{sec:ns} proves Theorem~\ref{thm:main-ns}. Section~\ref{sec:partI} proves
$R_3 \le 11$, Section~\ref{sec:partII} sets up the eleven-variable analysis and
Appendix~\ref{app:cases} carries it out in full, completing Theorem~\ref{thm:main-r3}.
Section~\ref{sec:comput} describes the computational verification. Section~\ref{sec:open}
collects open problems.

\section{Preliminaries}\label{sec:prelim}

\subsection{Fourier analysis on the cube}
We use the $\pm 1$ convention throughout: a Boolean function is a map
$f \colon \cube^n \to \cube$, and every $g \colon \cube^n \to \bbR$ has a unique multilinear
expansion $g = \sum_{S \subseteq [n]} \Fhat{g}(S) \chi_S$ with $\chi_S(x) = \prod_{i \in S}x_i$.
We identify a subset $S \subseteq [n]$ with its indicator vector in $\bbF_2^n$, so that
$\chi_S \chi_T = \chi_{S \sd T}$, where $\sd$ is symmetric difference, i.e.\ addition in
$\bbF_2^n$. Expectations $\Ex$ and probabilities $\Prb$ are with respect to the uniform measure
on $\cube^n$, and Parseval reads $\Ex[g^2] = \sum_S \Fhat{g}(S)^2$; in particular
$\sum_S \Fhat{f}(S)^2 = 1$ for Boolean $f$.

\subsection{Granularity}
\begin{fact}[Granularity; {\cite[Ex.~1.11]{ODbook}}, \cite{NS94}]\label{fact:gran}
If $f \colon \cube^n \to \cube$ has $\Deg f \le d$, then every Fourier coefficient
$\Fhat{f}(S)$ is an integer multiple of $2^{1-d}$.
\end{fact}
\todo{confirm the exercise number in O'Donnell's book; the statement itself is standard and is
already implicit in \cite{NS94}.}

Granularity is what makes the whole argument finite: for $\Deg f \le d$ the rescaled
coefficients $n_S := 2^{d-1}\Fhat{f}(S)$ are integers with $\sum_S n_S^2 = 4^{d-1}$, and the
problem becomes one about integer vectors.

\subsection{Derivatives, influences, links}
For $i \in [n]$ put
\[
  D_i f(x) := \frac{f(x^{i \to 1}) - f(x^{i \to -1})}{2}
  = \sum_{S \ni i} \Fhat{f}(S)\,\chi_{S \setminus i},
\]
where $x^{i \to b}$ is $x$ with the $i$-th coordinate set to $b$. If $f$ is Boolean then $D_i f$
is $\bits$-valued, and $\Deg D_i f \le \Deg f - 1$. The \emph{influence} of $i$ is
\[
  \Inf_i(f) = \Prb\left[ f(x) \ne f(x^{\oplus i}) \right] = \Prb[D_i f \ne 0]
  = \sum_{S \ni i} \Fhat{f}(S)^2 ,
\]
and the variable $x_i$ is relevant exactly when $\Inf_i(f) > 0$.

\begin{fact}[Support bound; {\cite{NS94}}]\label{fact:support}
A nonzero function $g \colon \cube^n \to \bbR$ of degree at most $k$ is nonzero on at least a
$2^{-k}$ fraction of the cube.
\end{fact}

\begin{proof}
Restrict to a subcube on which a top-degree monomial of $g$ survives and induct on $k$; this is
the Schwartz--Zippel-type bound used in \cite{NS94}.
\end{proof}

\begin{corollary}\label{cor:minfluence}
If $\Deg f \le d$ and $x_i$ is relevant, then $\Inf_i(f) \ge 2^{1-d}$.
\end{corollary}

\begin{proof}
$D_i f$ is nonzero, $\bits$-valued and of degree at most $d - 1$, so by Fact~\ref{fact:support}
$\Prb[D_i f \ne 0] \ge 2^{1-d}$.
\end{proof}

\begin{fact}[Total influence]\label{fact:totalinf}
$\sum_{i=1}^n \Inf_i(f) = \sum_S |S|\,\Fhat{f}(S)^2 \le \Deg f$ for Boolean $f$.
\end{fact}

Corollary~\ref{cor:minfluence} and Fact~\ref{fact:totalinf} give \eqref{eq:ns} at once:
$d \ge \sum_i \Inf_i \ge n\,2^{1-d}$, so $n \le d\,2^{d-1}$. Everything below is an analysis of
the near-equality case.

\begin{definition}[Mass, support hypergraph, link]\label{def:link}
Let $f \colon \cube^n \to \cube$ have $\Deg f \le d$ and put $n_S := 2^{d-1}\Fhat{f}(S)
\in \bbZ$. The \emph{support} of $f$ is $\mathcal{F} := \{S : n_S \ne 0\}$, regarded as a
hypergraph on the vertex set $[n]$; when all its members have size $d$ we call it
\emph{$d$-uniform}. For a vertex $i$ we write
\[
  m_i := \sum_{S \ni i} n_S^2 = 4^{d-1}\,\Inf_i(f)
\]
for the \emph{mass at $i$}, and we call the signed family
\[
  \lk(i) := \left\{\, (S \setminus i,\, n_S) \;:\; S \in \mathcal{F},\ S \ni i \,\right\}
\]
the \emph{link} of $i$. We often abuse notation and call the underlying set family
$\{S \setminus i : S \in \mathcal{F},\, S \ni i\}$ the link as well; when $\mathcal{F}$ is
$d$-uniform the link is a family of $(d-1)$-sets and, for $d = 3$, a graph.
\end{definition}

By definition $2^{d-1}D_i f = \sum_{S \ni i} n_S \chi_{S \setminus i}$ is
$\{0, \pm 2^{d-1}\}$-valued with $\sum_{S \ni i} n_S^2 = m_i$. The entire paper consists of
extracting structure from that one sentence.

\subsection{The Chiarelli--Hatami--Saks lower bound}
For completeness we record the construction from \cite[\S3]{CHS20} at $d = 2, 3$. Put
\[
  \XiC_2(a,b,u,w) := \frac{a+b}{2}\,u + \frac{a-b}{2}\,w
  = \frac{au + bu + aw - bw}{2},
\]
which is the degree-$2$ Boolean function \eqref{eq:degree2extremal} on $4$ variables, and
\[
  \XiC_3(s,t,x,y) := \frac{s+t}{2}\,\XiC_2(x) + \frac{s-t}{2}\,\XiC_2(y),
\]
where $x = (x_a,x_b,x_u,x_w)$ and $y = (y_a,y_b,y_u,y_w)$ are disjoint $4$-tuples of variables.
Expanding,
\[
  \XiC_3 = \frac14\Big[ (s+t)\big(x_ax_u + x_bx_u + x_ax_w - x_bx_w\big)
        + (s-t)\big(y_ay_u + y_by_u + y_ay_w - y_by_w\big) \Big],
\]
a homogeneous cubic with $16$ terms of coefficient $\pm\frac14$ on
$2 + 4 + 4 = 10$ relevant variables. Thus $R_3 \ge 10$. In general
$\XiC_d$ has $3 \cdot 2^{d-1} - 2$ relevant variables, giving \eqref{eq:chslower}.

\section{The Nisan--Szegedy bound is never attained for \texorpdfstring{$d \ge 3$}{d >= 3}}
\label{sec:ns}

Throughout this section $f \colon \cube^n \to \cube$ has $\Deg f \le d$ with all $n$ variables
relevant, and we write
\[
  m := d - 1, \qquad N := 2^m = 2^{d-1}.
\]

\subsection{Minimal mass forces an affine cube of unit terms}

\begin{lemma}[Minimal support forces an affine cube of unit terms]\label{lem:A}
Let
\[
  q = 2^{-m} \sum_{j=1}^{J} n_j \chi_{T_j}
\]
with the $T_j \subseteq [n]$ distinct, the $n_j$ nonzero integers, and $\sum_{j} n_j^2 = N$.
Suppose $q$ is $\bits$-valued and not identically $0$. Then all $n_j = \pm 1$, $J = N$, the
family $\{T_j\}$ is an affine subspace of dimension $m$ of $\bbF_2^n$, and
\[
  n_j n_k n_l n_r = +1 \qquad \text{whenever} \qquad T_j \sd T_k \sd T_l \sd T_r = \emptyset .
\]
\end{lemma}

\begin{proof}
Put $s := 2^m q = \sum_j n_j \chi_{T_j}$, so $s$ takes values in $\{0, \pm N\}$ and, by
Parseval, $\Ex[s^2] = \sum_j n_j^2 = N$; hence $\Prb[s \ne 0] = 1/N$ and
$\Ex[s^4] = N^2\,\Ex[s^2] = N^3$. Expanding the fourth power and using
$\Ex[\chi_{T_j}\chi_{T_k}\chi_{T_l}\chi_{T_r}] = \mathbf{1}\{T_j \sd T_k \sd T_l \sd T_r =
\emptyset\}$,
\[
  \Ex[s^4] = \sum_{\substack{(j,k,l,r) \\ T_j \sd T_k \sd T_l \sd T_r = \emptyset}}
  n_j n_k n_l n_r .
\]
For each ordered triple $(j,k,l)$ there is \emph{at most one} index $r$ with
$T_r = T_j \sd T_k \sd T_l$, because the $T$'s are distinct; call the set of triples possessing
such an $r$ the \emph{dominant} set $\mathrm{Dom}$, and write $r(j,k,l)$ for that index. Then
\begin{align}
  N^3 = \Ex[s^4]
  &\le \sum_{(j,k,l) \in \mathrm{Dom}} |n_j n_k n_l|\,|n_{r(j,k,l)}| \notag \\
  &\le \sum_{(j,k,l) \in \mathrm{Dom}} n_j^2 n_k^2\, |n_l|\,|n_{r(j,k,l)}|
    \label{eq:Astep2} \\
  &\le \Big( \sum_{\mathrm{Dom}} n_j^2 n_k^2 n_l^2 \Big)^{1/2}
       \Big( \sum_{\mathrm{Dom}} n_j^2 n_k^2 n_{r(j,k,l)}^2 \Big)^{1/2}
    \label{eq:ACS} \\
  &\le \big( N^3 \big)^{1/2}\big( N^3 \big)^{1/2} = N^3 . \notag
\end{align}
Line \eqref{eq:Astep2} uses $|n| \le n^2$ for nonzero integers $n$, applied to the factors
$n_j$ and $n_k$; line \eqref{eq:ACS} is Cauchy--Schwarz. The final bound uses
$\sum_{\mathrm{Dom}} n_j^2 n_k^2 n_l^2 \le \big(\sum_j n_j^2\big)^3 = N^3$ and, for the second
factor, that for fixed $(j,k)$ the map $l \mapsto r(j,k,l)$ is injective on $\mathrm{Dom}$
(since $l$ is recovered from $r$ by $T_l = T_j \sd T_k \sd T_r$), so that
$\sum_l n_{r(j,k,l)}^2 \le \sum_j n_j^2 = N$ and the whole sum is at most
$\big(\sum n_j^2\big)^2 \cdot N = N^3$.

Equality throughout forces the following.
\begin{enumerate}[label=(\roman*)]
  \item $\mathrm{Dom}$ is the set of \emph{all} triples: the family $\{T_j\}$ is closed under
  $(x,y,z) \mapsto x \sd y \sd z$, hence is an affine subspace of $\bbF_2^n$, of some dimension
  $k$, and $J = 2^k$.
  \item Equality in \eqref{eq:Astep2} forces $|n_j n_k| = 1$ whenever the triple contributes,
  and equality in Cauchy--Schwarz gives $|n_{r(j,k,l)}| = \lambda |n_l|$ for all $j,k,l$ and
  some $\lambda > 0$; taking $j = k$ gives $r = l$ and $\lambda = 1$, and for fixed $k,l$ the
  map $j \mapsto r(j,k,l)$ is onto, so all $|n_j|$ equal a common value $c$. (In fact $c = 1$
  already follows from \eqref{eq:Astep2}; the argument is kept because it also delivers the
  count $J$.)
  \item All products $n_j n_k n_l n_{r(j,k,l)}$ are positive.
\end{enumerate}
From (ii), $2^k c^2 = N = 2^m$, so $c = 2^t$ with $k = m - 2t$. If $t \ge 1$ then
$|s| \le c\,2^k = 2^{m-t} < N$ pointwise, so $s \in \{0,\pm N\}$ forces $s \equiv 0$,
contradicting $q \ne 0$. Hence $c = 1$, $k = m$ and $J = N$.
\end{proof}

\begin{remark}\label{rem:A}
Item (iii) says that $\eps_j := n_j$ is an affine character on the affine subspace $\{T_j\}$:
writing $\{T_j\} = T_0 \sd L$ for a linear subspace $L$ of dimension $m$, we have
$\eps_{T_0 \sd A} = \eps_0\, \psi(A)$ for a homomorphism $\psi \colon L \to \{\pm 1\}$.
Consequently, for any basis $D_1, \dots, D_m$ of $L$,
\[
  q = \pm\, \chi_{T_0} \prod_{i=1}^{m} \frac{1 + \psi(D_i)\,\chi_{D_i}}{2},
\]
that is: a minimal-mass derivative is a character times the indicator of an affine subspace of
codimension $m = d - 1$. The converse also holds, every such product being $\bits$-valued.
Note that Lemma~\ref{lem:A} uses no degree hypothesis at all; it is a statement about sparse,
granular, $\bits$-valued functions, and it applies to any variable of minimal influence
$2^{1-d}$ in any degree-$d$ Boolean function, homogeneous or not.
\end{remark}

\subsection{Equal sizes force a cross-polytope}

\begin{lemma}[Equal sizes force a cross-polytope]\label{lem:B}
Let $\{T_p\} = T_0 \sd L$ be an affine subspace of dimension $m$ of subsets of $[n]$, all of
whose members have size exactly $m$. Then there are $m$ pairwise disjoint pairs
$\{t_i, t_i'\}$, $i = 1, \dots, m$, with $T_0 = \{t_1, \dots, t_m\}$ and
\[
  \{T_p\} = \big\{\, \{u_1, \dots, u_m\} \;:\; u_i \in \{t_i, t_i'\} \,\big\},
\]
that is, $\{T_p\}$ is the facet family of an $m$-dimensional cross-polytope on the $2m$
vertices $t_1, t_1', \dots, t_m, t_m'$.
\end{lemma}

\begin{proof}
For $A \in L \setminus \{\emptyset\}$ and any $p$, the equality $|T_p \sd A| = m = |T_p|$ gives
$|A| = 2\,|A \cap T_p|$. Applying this with $T_p = T_0 \sd B$ for $B \in L$,
\[
  \frac{|A|}{2} = |(T_0 \sd B) \cap A| = |T_0 \cap A| + |B \cap A| - 2\,|T_0 \cap A \cap B| ,
\]
and since $|A|/2 = |T_0 \cap A|$ we obtain
\[
  |A \cap B| = 2\,|T_0 \cap A \cap B| \qquad \text{for all } A, B \in L.
\]
The $\bbF_2$-linear map $L \to 2^{T_0}$, $A \mapsto A \cap T_0$, is injective (if
$A \cap T_0 = \emptyset$ then $|A| = 2|A \cap T_0| = 0$), and both spaces have dimension $m$, so
it is a bijection. For $t \in T_0$ let $A_t \in L$ be the element with $A_t \cap T_0 = \{t\}$;
then $|A_t| = 2$, say $A_t = \{t, t'\}$ with $t' \notin T_0$. For $t \ne s$ in $T_0$,
$|A_t \cap A_s| = 2\,|T_0 \cap A_t \cap A_s| = 0$, so the elements $t'$ are distinct and the
pairs $\{t,t'\}$ are disjoint. The $A_t$ form a basis of $L$, and for $P \subseteq T_0$,
\[
  T_0 \sd \Big( \mathop{\triangle}\limits_{t \in P} A_t \Big)
  = (T_0 \setminus P) \cup \{ t' : t \in P \},
\]
which is exactly the asserted description.
\end{proof}

\begin{fact}[Antipodal pairs are readable from the facets]\label{fact:antipodal}
The facet family of an $m$-dimensional cross-polytope determines its antipodal pairs: two
vertices are antipodal if and only if they never occur in a common facet, every non-antipodal
pair lying in a common facet.
\end{fact}

We use Fact~\ref{fact:antipodal} to read off the pairs of a link from its facets.

\subsection{Disjoint sums}

\begin{lemma}[Disjoint sums]\label{lem:C}
Let $g$ and $h$ be real functions of disjoint sets of variables with $g(x) + h(y) \in \{\pm1\}$
for all $x, y$. If $g$ is not constant then $h$ is constant.
\end{lemma}

\begin{proof}
Let $u \ne u'$ be two values of $g$. Then $h(y) \in \{1-u, -1-u\} \cap \{1-u', -1-u'\}$ for
every $y$, and two distinct two-element sets meet in at most one point, so $h$ is constant.
\end{proof}

In particular a sum $f_1(x) + f_2(y)$ of two non-constant functions of disjoint variable sets is
never $\pm1$-valued, since by Lemma~\ref{lem:C} one of them must be constant. We use this
repeatedly, always in the following form: if the support of a Boolean $f$ splits into two
nonempty parts on disjoint variable sets, each carrying positive Fourier weight on
non-constant characters, we have a contradiction.

\subsection{Proof of Theorem~\ref{thm:main-ns}}

\begin{proof}[Proof of Theorem~\ref{thm:main-ns}]
Let $d \ge 3$ and suppose $f \colon \cube^n \to \cube$ has $\Deg f \le d$ and
$n = d\,2^{d-1}$ relevant variables.

\smallskip\noindent\textbf{Step 1: the extremal shape.}
By Corollary~\ref{cor:minfluence} and Fact~\ref{fact:totalinf},
\[
  d \;\ge\; \Deg f \;\ge\; \sum_{i} \Inf_i(f) \;\ge\; n\,2^{1-d} = d,
\]
so equality holds everywhere: $\Deg f = d$ exactly, every $\Inf_i(f) = 2^{1-d}$, and
$\sum_S |S|\Fhat{f}(S)^2 = d = d\sum_S \Fhat{f}(S)^2$, which forces all Fourier weight onto
level $d$. Thus $f$ is a homogeneous degree-$d$ multilinear polynomial,
$f = \sum_{|S| = d} c_S \chi_S$.

\smallskip\noindent\textbf{Step 2: every link is a cross-polytope.}
Fix a relevant variable $i$ and write $D_i f = 2^{-m}\sum_{S \ni i} n_S \chi_{S \setminus i}$
with $n_S := 2^m c_S \in \bbZ$ by Fact~\ref{fact:gran}. Then
$\sum_{S \ni i} n_S^2 = 2^{2m}\Inf_i(f) = 2^m = N$, the function $D_i f$ is $\bits$-valued and
nonzero, and the sets $S \setminus i$ (with $i$ fixed) are distinct. Lemma~\ref{lem:A} applies
and gives: every $n_S$ with $S \ni i$ equals $\pm 1$, so every coefficient of $f$ is
$\pm 2^{-m}$ and, by Parseval, $f$ has exactly $2^{2m}$ terms, forming a $d$-uniform hypergraph
$\mathcal{F}$ in which every vertex lies in exactly $N = 2^{d-1}$ terms; and the link
$\lk(i) = \{S \setminus i : S \in \mathcal{F},\, S \ni i\}$ is an affine subspace of dimension
$m$ consisting of $m$-sets. By Lemma~\ref{lem:B}, $\lk(i)$ is the facet family of an
$m$-dimensional cross-polytope: there are $m$ disjoint pairs
$P_1(i), \dots, P_m(i)$ of vertices (the \emph{antipodal pairs at $i$}), and the terms at $i$
are exactly the sets $\{i\} \cup \{ \text{one vertex from each } P_k(i)\}$. In particular two
vertices antipodal at $i$ never lie in a common term with $i$, and $i$ has exactly $2m$
neighbours.

\smallskip\noindent\textbf{Step 3: closure.}
Fix a vertex $v$ with pairs $P_k = \{t_k, t_k'\}$, $k = 1, \dots, m$; note $m \ge 2$ because
$d \ge 3$. Put $a := t_1$. The terms at $v$ containing $a$ are the sets
$\{v,a\} \cup (\text{one from each } P_2, \dots, P_m)$; removing $a$, these are $2^{m-1}$ facets
of the cross-polytope $\lk(a)$, all containing $v$. By Fact~\ref{fact:antipodal} the pairs of
$\lk(a)$ are therefore $P_2, \dots, P_m$ together with one further pair $\{v, w\}$ for a unique
vertex $w$, the \emph{antipode of $v$ at $a$}. Clearly
$w \notin \{v\} \cup P_2 \cup \dots \cup P_m$. Also $w \ne t_1'$: otherwise the term
$\{a, t_1', t_2 \text{ or } t_2', \dots\}$ exists, and removing $t_2$ from it exhibits a facet
of $\lk(t_2)$ containing both $a = t_1$ and $t_1'$ --- but $P_1$ is a pair of $\lk(t_2)$ (as the
terms at $v$ containing $t_2$ show), so $t_1$ and $t_1'$ are antipodal at $t_2$, and
Fact~\ref{fact:antipodal} forbids their lying in a common facet. So $w$ is a new vertex.

The terms at $a$ are therefore $\{a\} \cup \{v \text{ or } w\} \cup
(\text{one from each } P_2, \dots, P_m)$. In particular
$\{a, w, t_2, u_3, \dots, u_m\}$ is a term for suitable $u_k \in P_k$, so removing $t_2$ gives a
facet of $\lk(t_2)$ containing $a$ and $w$. The pairs of $\lk(t_2)$ are
$P_1, P_3, \dots, P_m$ together with $\{v, w_2\}$, where $w_2$ is the antipode of $v$ at $t_2$;
since $w$ appears in a facet of $\lk(t_2)$ and is none of the vertices of
$P_1, P_3, \dots, P_m$, we get $w = w_2$. Symmetrically the antipode of $v$ at every $t_k$ and
every $t_k'$ equals $w$. Hence $w$ is adjacent to all $2m$ vertices of
$P_1 \cup \dots \cup P_m$; these exhaust the $2m$ neighbours of $w$, and the terms at $w$ are
exactly $\{w\} \cup (\text{one from each } P_k)$.

Therefore the connected component of $v$ in $\mathcal{F}$ is exactly the boundary complex of the
$d$-dimensional cross-polytope on the $2d$ vertices $\{v, w\} \cup P_1 \cup \dots \cup P_m$:
it has $2d$ vertices and $2^d$ facets, and no other term of $f$ touches these vertices, since
every one of them already lies in its full quota of $2^{d-1}$ terms.

\smallskip\noindent\textbf{Step 4: at least two components, and the contradiction.}
By Step 2, $\mathcal{F}$ has $2^{2m} = 4^{d-1}$ terms, and by Step 3 it is a disjoint union of
cross-polytope boundaries, each with $2^d$ terms on its own $2d$ vertices. Hence there are
exactly
\[
  \frac{2^{2m}}{2^{d}} = 2^{d-2} \;\ge\; 2 \qquad (d \ge 3)
\]
components, on pairwise disjoint variable sets. Writing $f = f_1 + g$, where $f_1$ collects the
$2^d$ terms of the first component (on its $2d$ variables) and $g$ the rest, both $f_1$ and $g$
are non-constant, since each carries positive Fourier weight on non-constant characters. This
contradicts Lemma~\ref{lem:C}. Therefore no degree-$d$ Boolean function has $d\,2^{d-1}$
relevant variables, and $R_d \le d\,2^{d-1} - 1$ for $d \ge 3$.
\end{proof}

\begin{remark}
The exponent $2^{d-2} \ge 2$ in Step 4 is exactly where $d \ge 3$ is used, and it is the only
place. For $d = 2$ one has $2^{d-2} = 1$: the single component is the boundary of the
$2$-dimensional cross-polytope, i.e.\ a square, and the four edges with signs $(+,+,+,-)$ give
the Boolean function \eqref{eq:degree2extremal} with $4 = 2 \cdot 2^1$ relevant variables.
Step 3 also needs $m \ge 2$, i.e.\ $d \ge 3$, for the vertex $t_2$ to exist. So the proof
correctly fails at $d = 2$, as it must.
\end{remark}

\section{Part I: no degree-3 Boolean function has 12 relevant variables}\label{sec:partI}

This section proves the case $d = 3$ of Theorem~\ref{thm:main-ns} in a self-contained form,
which is also the first half of Theorem~\ref{thm:main-r3}. The reader who has read
Section~\ref{sec:ns} may skip to Section~\ref{sec:partII}; we keep the $d=3$ argument because it
runs through a closed-surface classification that is used again, in a more delicate form, in
Part~II.

\begin{theorem}\label{thm:no12}
No Boolean function of degree $3$ has $12$ relevant variables. Hence $R_3 \le 11$.
\end{theorem}

\begin{proof}
\textbf{Step 1: twelve relevant variables force a homogeneous cubic with $16$ terms $\pm\frac14$,
$4$-regular.}
With $12$ relevant variables, Corollary~\ref{cor:minfluence} and Fact~\ref{fact:totalinf} give
$3 \ge \sum_i \Inf_i \ge 12 \cdot \frac14 = 3$. Hence every $\Inf_i = \frac14$ exactly, and
$\sum_S |S| \Fhat{f}(S)^2 = 3 = 3\sum_S \Fhat{f}(S)^2$, so all Fourier weight sits on level
$3$: $f = \sum_{|S| = 3} c_S \chi_S$. By Fact~\ref{fact:gran}, $n_S := 4 c_S$ is an integer,
$\sum_S n_S^2 = 16$ by Parseval, and $\Inf_i = \sum_{S \ni i} c_S^2 = \frac14$ gives
$m_i = \sum_{S \ni i} n_S^2 = 4$ for every vertex $i$. So at each vertex there is either one
term with $|n_S| = 2$, or four terms with $|n_S| = 1$.

Suppose some $S$ had $|n_S| = 2$. Then the three vertices of $S$ lie in no other term, so
$f = \pm\frac12 \chi_S + g$ with $g$ a function of the other $9$ variables. Then
$f^2 = \frac14 + g^2 \pm \chi_S g = 1$; the monomials of $\chi_S g$ all contain $S$ while those
of $\frac34 - g^2$ avoid $S$, so $\chi_S g = 0$, i.e.\ $g = 0$, contradicting
$\sum \Fhat{g}(\cdot)^2 = \frac34$. Hence all $n_S = \pm 1$:
\[
  f = \frac14 \sum_{S \in \mathcal{F}} \eps_S \chi_S,
\]
with $\mathcal{F}$ a $3$-uniform hypergraph on the $12$ vertices, $|\mathcal{F}| = 16$, and
every vertex in exactly $4$ triples.

\smallskip\noindent\textbf{Step 2: every vertex link is a $4$-cycle.}
Fix a vertex $i$. Then
$D_i f = \sum_{S \ni i} c_S \chi_{S \setminus i} = \frac14\sum_{j=1}^4 \eps_j \chi_{T_j}$, where
$T_j = S_j \setminus \{i\}$ are four distinct pairs. Since $D_i f$ is $\bits$-valued, writing
$u_j = \eps_j \chi_{T_j}(x) \in \{\pm1\}$ we need $\sum_j u_j \in \{0, \pm 4\}$ for every $x$,
i.e.\ never $\pm 2$, i.e.\ the number of $j$ with $u_j = -1$ is always even, i.e.\
$\prod_j u_j \equiv 1$. That is,
$\big(\prod_j \eps_j\big)\chi_{T_1 \sd T_2 \sd T_3 \sd T_4} \equiv 1$, so
\[
  T_1 \sd T_2 \sd T_3 \sd T_4 = \emptyset \qquad \text{and} \qquad \prod_j \eps_j = +1 .
\]
Four distinct pairs with empty symmetric difference: every vertex lies in an even number of
them, there are $8$ incidences, and a vertex lying in all four would leave the other four
endpoints in one pair each (impossible, as those pairs would then have odd degrees). So exactly
four vertices lie in exactly two pairs each, i.e.\ the pairs form the unique $2$-regular simple
graph on $4$ vertices, a $4$-cycle $a$--$b$--$c$--$d$--$a$. Consequently the four triples at $i$
are $\{i,a,b\}, \{i,b,c\}, \{i,c,d\}, \{i,d,a\}$: the link of $i$ in $\mathcal{F}$ is $C_4$, and
every pair $\{i,x\}$ lies in exactly $2$ triples (for $x \in \{a,b,c,d\}$) or in none.

\smallskip\noindent\textbf{Step 3: $\mathcal{F}$ is a closed triangulated surface with all
vertex degrees $4$, hence two disjoint octahedra.}
By Step 2 the pure $2$-dimensional complex $|\mathcal{F}|$ has every edge in exactly two
triangles and every vertex link a single cycle $C_4$, so it is a closed $2$-manifold (possibly
disconnected), triangulated with $V = 12$ vertices, $F = 16$ faces and
$E = (3 \cdot 16)/2 = 24$ edges. Its Euler characteristic is $V - E + F = 12 - 24 + 16 = 4$. A
connected closed surface has $\chi \le 2$, so there are at least two components. Each component
is a closed surface all of whose vertices have degree $4$, so with $V'$ vertices it has
$E' = 2V'$ and $F' = 4V'/3$, whence $\chi' = V' - 2V' + 4V'/3 = V'/3 \in \{1, 2\}$. The value
$\chi' = 1$ would need $V' = 3$, impossible for a complex with a degree-$4$ vertex (and a
triangulated projective plane needs at least $6$ vertices in any case). So $\chi' = 2$ and
$V' = 6$: the component is the octahedron, which is the unique $4$-regular triangulated sphere
(the complement of a $4$-regular graph on $6$ vertices is a perfect matching, so the graph is
$K_{2,2,2}$ and all $8$ faces are then forced).

There is also a topology-free route to the same conclusion: if $\lk(i) = a$--$b$--$c$--$d$ then
$\lk(a) = b$--$i$--$d$--$x$ for some $x \ne c$ (were $x = c$, the link of $b$ would contain a
$3$-cycle, impossible inside a $C_4$), and the closure forces the octahedron on
$\{i,a,b,c,d,x\}$; this was checked exhaustively by an independent referee.

Two octahedra use all $12$ vertices. Hence $\mathcal{F}$ is the set of faces of two
vertex-disjoint octahedra and
\[
  f = f_1(x_1,\dots,x_6) + f_2(x_7,\dots,x_{12}),
\]
with $f_1, f_2$ each carrying $8$ terms $\pm\frac14$ on disjoint variable sets, neither
constant.

\smallskip\noindent\textbf{Step 4: a sum of two non-constant functions on disjoint variable sets
is never Boolean.}
This is Lemma~\ref{lem:C}: if $f_1(x) + f_2(y) \in \{\pm1\}$ for all $x, y$ and $f_1$ takes two
values $u \ne u'$, then $f_2(y) \in \{1-u,-1-u\} \cap \{1-u',-1-u'\}$ for all $y$, an
intersection with at most one element, so $f_2$ is constant --- contradicting that $f_2$ has
Fourier weight $\frac12$ on non-constant characters.

Equivalently and more directly: for disjoint triples $S$ in the first octahedron and $T$ in the
second, the coefficient of $\chi_{S \cup T}$ in $f^2$ is $2\eps_S\eps_T/16 \ne 0$, since no
other pair of triples of $\mathcal{F}$ has union $S \cup T$; but $f^2 \equiv 1$, so every
non-constant coefficient of $f^2$ vanishes. Contradiction.
\end{proof}

\section{Part II: the eleven-variable case}\label{sec:partII}

Throughout this section $f \colon \cube^{11} \to \cube$ has $\Deg f \le 3$ and all $11$
variables relevant; we derive a contradiction. Write $f = \sum_S c_S \chi_S$ and
$n_S := 4 c_S \in \bbZ$ by Fact~\ref{fact:gran}, so that $\sum_S n_S^2 = 16$ by Parseval, and
for a vertex $i$ put
\[
  m_i := \sum_{S \ni i} n_S^2 = 16\,\Inf_i(f) = 16\,\Prb[D_i f \ne 0].
\]

Compared with Part~I there is one unit of slack in the influence budget: eleven variables need
total influence at least $11/4$, while $3$ is available. The two lemmas below say exactly how
that slack can be spent, and the bookkeeping identity of Section~\ref{sec:bookkeep} turns it
into a finite list of configurations, which Appendix~\ref{app:cases} eliminates.

\subsection{The masses}

\begin{lemma}[Values of $m_i$]\label{lem:masses}
$m_i \in \{4, 6, 8\}$ for every $i$, and $\sum_i m_i = \sum_S |S|\,n_S^2 \le 48$.
\end{lemma}

\begin{proof}
The derivative $D_i f = \sum_{S \ni i} (n_S/4)\chi_{S \setminus i}$ is $\bits$-valued (because
$f$ is Boolean) and of degree at most $2$, so
$4 D_i f = \sum_j n_j \chi_{T_j}$ takes values in $\{0, \pm 4\}$, where the $T_j$ are distinct
sets of size at most $2$ and $\sum_j n_j^2 = m_i$. A relevant variable has
$\Inf_i \ge \frac14$ by Corollary~\ref{cor:minfluence}, so $m_i \ge 4$. Also
$\sum_i m_i = \sum_S |S| n_S^2 \le 3\sum_S n_S^2 = 48$, so with $11$ vertices each of mass at
least $4$ we get $m_i \le 48 - 40 = 8$ for every $i$.

It remains to exclude $m_i = 5$ and $m_i = 7$.
\begin{itemize}[leftmargin=1.4em]
  \item Any pattern with an odd number of odd $n_j$ has odd value $\sum_j n_j \chi_{T_j}(x)$ at
  every point, hence never lies in $\{0, \pm4\}$. This kills five $\pm1$'s, seven $\pm1$'s,
  $\{\pm2,\pm1\}$ and $\{\pm2,\pm1,\pm1,\pm1\}$.
  \item The pattern $\{\pm2,\pm1,\pm1\}$ (mass $6$): $2u_0 + u_1 + u_2 \in \{0,\pm4\}$ with
  $u_k \in \{\pm1\}$ forces $u_1 = u_2$ for every $x$, i.e.\
  $\chi_{T_1} \equiv \pm \chi_{T_2}$, impossible for distinct $T_1 \ne T_2$. So mass $6$ occurs
  only as six terms $\pm1$.
\end{itemize}
Since $m_i \le 8$ forces $|n_j| \le 2$ throughout, the listed patterns are all the ones that
occur, and $m_i \in \{4,6,8\}$.
\end{proof}

\begin{remark}
A shorter route to the parity part: $4D_i f$ is $\{0,\pm4\}$-valued and
$\sum_j n_j \equiv \sum_j n_j^2 = m_i \pmod 2$, so $m_i$ is even.
\end{remark}

\subsection{The link shapes}

\begin{lemma}[Link structure]\label{lem:links}
With the link of $i$ as in Definition~\ref{def:link}, written
$\{T_j\} = \{S \setminus i : n_S \ne 0,\, S \ni i\}$ with signs $\eps_j$:
\begin{enumerate}[label=\textup{(\alph*)}]
  \item \textbf{$m_i = 4$}: four terms $\pm1$; the condition is necessary and sufficient that
  $T_1 \sd T_2 \sd T_3 \sd T_4 = \emptyset$ and $\prod_j \eps_j = +1$ (a sum of four $\pm1$'s
  avoids $\pm2$ iff the number of $-1$'s is even iff $\prod_j u_j \equiv 1$). Four distinct
  sets of size at most $2$ in which every element lies in an even number of members are exactly:
  \begin{itemize}[leftmargin=1.4em]
    \item[(L1)] a $4$-cycle of pairs on $4$ vertices;
    \item[(L2)] $\{\emptyset, \{a\}, \{b\}, \{a,b\}\}$;
    \item[(L3)] $\{\{a\}, \{b\}, \{a,c\}, \{b,c\}\}$;
    \item[(L4)] $\{\emptyset, \{a,b\}, \{b,c\}, \{a,c\}\}$.
  \end{itemize}
  (Enumeration, verified by computer: without $\emptyset$, four pairs with all degrees even form
  a $C_4$; singletons come in pairs $\{a\},\{b\}$ needing $\{a,b\}$ plus a fourth set ---
  $\emptyset$ gives (L2) --- or $\{a,c\},\{b,c\}$ giving (L3); with $\emptyset$ present the other
  three sets have symmetric difference $\emptyset$, so they are three pairs forming a triangle,
  giving (L4), or $\{a\},\{b\},\{a,b\}$, giving (L2) again.)
  In terms of $f$: (L1) means $i$ lies in four cubic terms only; (L2) means a linear term
  $\{i\}$, quadratics $\{i,a\},\{i,b\}$ and a cubic $\{i,a,b\}$; (L3) means quadratics
  $\{i,a\},\{i,b\}$ and cubics $\{i,a,c\},\{i,b,c\}$; (L4) means a linear term $\{i\}$ and
  cubics $\{i,a,b\},\{i,b,c\},\{i,a,c\}$. In particular: a link containing a singleton is (L2)
  or (L3); a link containing $\emptyset$ is (L2) or (L4); a link containing a triangle of pairs
  is (L4), a $C_4$ having no triangle.
  \item \textbf{$m_i = 6$}: six terms $\pm1$ with $T_1 \sd \dots \sd T_6 = \emptyset$ and
  $\prod_j \eps_j = -1$ (six $\pm1$'s avoid $\pm2$ and $\pm6$ iff the number of $-1$'s is odd).
  \item \textbf{$m_i = 8$}: either two terms $\pm2$ (on any two distinct sets), or one $\pm2$ and
  four $\pm1$'s whose sets have symmetric difference $\emptyset$ and sign product $-1$, or eight
  terms $\pm1$ with symmetric difference $\emptyset$, sign product $+1$, and never all eight
  values agreeing.
  \item A vertex whose only term has $|n_S| = 2$ is impossible, since then
  $D_i f = \pm\frac12\chi_{S \setminus i}$ is not $\bits$-valued. Consequently a term with
  $|n_S| = 2$ forces every vertex of $S$ to have mass $8$ (mass $6$ cannot contain a $\pm2$, by
  the proof of Lemma~\ref{lem:masses}).
  \item In a link consisting of pairs only (no lower-order term at $i$), the multiplicity
  $\lambda_{ix} := \#\{S \in \mathcal{F} : S \ni i, x\}$ equals the degree of $x$ in the link
  graph; for (L1) it is $2$ for the four link vertices and $0$ otherwise.
\end{enumerate}
\end{lemma}

\begin{proof}
Each part is the value-set analysis of $4D_i f = \sum_j n_j \chi_{T_j}$ carried out as in
Lemma~\ref{lem:masses} and Step 2 of Theorem~\ref{thm:no12}: the sum of $k$ terms $\pm1$ lies in
$\{0,\pm4\}$ at every point iff the symmetric difference of the $T_j$ is empty with the stated
sign product, and the enumeration of admissible set families of size $4$ is the elementary
parity argument displayed in (a). Part (d) is immediate, and its consequence follows from
Lemma~\ref{lem:masses}. Part (e) is a restatement of the definition.
\end{proof}

\subsection{Bookkeeping}\label{sec:bookkeep}
Put
\[
  \delta := \sum_S (3 - |S|)\, n_S^2 \qquad\text{(the Fourier weight below the top level,
  weighted by codegree)}, \qquad e := \sum_{i} (m_i - 4).
\]
Then $\sum_i m_i = \sum_S |S| n_S^2 = 48 - \delta$ and $\sum_i m_i = 44 + e$, so
\begin{equation}\label{eq:bookkeeping}
  e + \delta = 48 - 44 = 4, \qquad e \ge 0, \quad \delta \ge 0 .
\end{equation}
The contributions to $\delta$ are: a quadratic term with $|n_S| = 1$ contributes $1$; a
quadratic with $|n_S| = 2$ contributes $4$; a linear term $\pm1$ contributes $2$; the constant
term $\pm1$ contributes $3$. A linear or constant term with $|n_S| = 2$ would contribute $8$ or
$12 > 4$ and is therefore impossible.

By Lemma~\ref{lem:masses}, $e \in \{0, 2, 4\}$, with $e = 0$ meaning all $m_i = 4$, $e = 2$
meaning exactly one vertex of mass $6$, and $e = 4$ meaning either two vertices of mass $6$ or
one of mass $8$. Correspondingly $\delta \in \{4, 2, 0\}$.

\subsection{A surface fact}

\begin{fact}[Surfaces of non-positive Euler characteristic]\label{fact:T}
In a closed connected triangulated surface with $V'$ vertices, $E'$ edges, $F'$ faces and Euler
characteristic $\chi'$, the relations $3F' = 2E'$ and $V' - E' + F' = \chi'$ give the degree sum
$2E' = 6V' - 6\chi'$. Hence $\chi' \le 0$ forces average degree at least $6$. In every complex
arising below, all degrees are $4$ except at most two vertices of degree at most $8$, so the
average degree is at most $4 + 8/V' < 6$ as soon as $V' \ge 5$; and $V' \ge 5$ holds because
some vertex has degree at least $4$ (a closed triangulated surface has no vertex of degree less
than $3$, and $V' \ge \deg + 1$). Therefore every component occurring below is a sphere (degree
sum $6V' - 12$) or a projective plane (degree sum $6V' - 6$, requiring at least $6$ vertices, by
Step~3 of Theorem~\ref{thm:no12}).
\end{fact}

Fact~\ref{fact:T} was checked by computer: for every degree pattern used in
Appendix~\ref{app:cases}, the relation $V' = 3\chi' + c/2$ (with $c$ the excess degree) was
tabulated, and every row with $\chi' \le 0$ gives $V' \le 2$.

\subsection{Statement and proof of Part II}

\begin{theorem}\label{thm:no11}
No Boolean function of degree $3$ has $11$ relevant variables.
\end{theorem}

\begin{proof}
By \eqref{eq:bookkeeping} we have $\delta \in \{4, 2, 0\}$. The three cases are eliminated in
Appendix~\ref{app:cases}: $\delta = 4$ in Section~\ref{app:delta4}, $\delta = 2$ in
Section~\ref{app:delta2}, and $\delta = 0$ in Section~\ref{app:delta0}. Since the three cases
are exhaustive, no such $f$ exists.
\end{proof}

\begin{proof}[Proof of Theorem~\ref{thm:main-r3}]
The function $\XiC_3$ of Section~\ref{sec:prelim} is Boolean of degree $3$ with $10$ relevant
variables, so $R_3 \ge 10$. Theorem~\ref{thm:no12} rules out $12$ relevant variables and
Theorem~\ref{thm:no11} rules out $11$. If some degree-$3$ Boolean function had
$n \ge 12$ relevant variables then $n \le 12$ by \eqref{eq:ns}, so $n = 12$, contradicting
Theorem~\ref{thm:no12}. Hence $R_3 = 10$.
\end{proof}

\section{Computational verification}\label{sec:comput}

The proofs above are complete as mathematics. Because the case analysis of
Appendix~\ref{app:cases} is long, we also verified the conclusion $R_3 = 10$ by three
independent computational routes: an exhaustive search over a finite reformulation, an
independent SAT-based structure search performed by a referee, and a partial formalisation in
Lean~4. This section states precisely what each covers.

\subsection{A finite reformulation}\label{sec:finite}
Let $n \ge 1$. A \emph{coefficient vector} is an integer-valued map
$N \colon \{S \subseteq [n] : |S| \le 3\} \to \bbZ$; write $n_S := N(S)$ and
\[
  f_N(x) := \frac14 \sum_{S} n_S \chi_S(x), \qquad x \in \cube^n .
\]

\begin{definition}[The statement $F(n)$]
$F(n)$ asserts that there is \emph{no} coefficient vector $N$ with
\begin{enumerate}[label=\textup{(\roman*)}]
  \item $\sum_S n_S^2 = 16$;
  \item for every nonempty $U \subseteq [n]$,
  $\displaystyle\sum_{\substack{(S,T)\ \text{ordered} \\ S \sd T = U}} n_S n_T = 0$;
  \item every $i \in [n]$ lies in some $S$ with $n_S \ne 0$.
\end{enumerate}
\end{definition}

Conditions (i) and (ii) together say exactly that $f_N^2 \equiv 1$, i.e.\ that $f_N$ is
$\cube$-valued of degree at most $3$: since $\chi_S\chi_T = \chi_{S \sd T}$, the expansion
$f_N^2 = \frac1{16}\sum_{S,T} n_S n_T \chi_{S \sd T}$ is identically $1$ iff its $U = \emptyset$
coefficient is $16$ and all others vanish. Condition (iii) says all $n$ variables are relevant
(a variable is relevant exactly when it appears in the Fourier support). Conversely, if
$f \colon \cube^n \to \cube$ has degree at most $3$ then Fact~\ref{fact:gran} makes
$N := 4\Fhat{f}$ a coefficient vector with $f = f_N$, and (i)--(iii) hold. Hence:
\[
  F(n) \ \Longrightarrow\ \text{no degree-3 Boolean function has } n \text{ relevant variables}.
\]
In particular $F(11)$ implies $R_3 \le 10$, and with $\XiC_3$ (which witnesses that $F(10)$ is
false) it gives $R_3 = 10$.

The search below is allowed to use only the following elementary consequences of $f_N^2 \equiv
1$, each provable in two lines and each independent of the hand proof:
\begin{itemize}[leftmargin=1.4em]
  \item[(E1)] For each $i$, $D_i f_N = \frac14\sum_{S \ni i} n_S \chi_{S \setminus i}$ is
  $\bits$-valued, so $4 D_i f_N$ takes values in $\{0, \pm4\}$ at every point.
  \item[(E2)] Parseval: $\sum_{S \ni i} n_S^2 = 16\,\Prb[D_i f_N \ne 0]$ is an integer in
  $[4,16]$ for relevant $i$, and $\sum_i \sum_{S \ni i} n_S^2 = \sum_S |S| n_S^2 \le 48$.
  \item[(E3)] Symmetries preserving (i)--(iii): permutations of $[n]$; negation of a variable
  $i$ (which flips $n_S$ for $S \ni i$); and the global sign $N \mapsto -N$.
\end{itemize}
Nothing beyond (E1)--(E3) --- in particular no link classification from
Section~\ref{sec:partII} --- is assumed by the search.

\subsection{Exhaustive canonical search}
The search program processes vertices one at a time; processing $v$ fixes every coefficient
$n_S$ with $v \in S$, that is, it chooses the link
$L_v = 4 D_v f_N = \sum_{S \ni v} n_S \chi_{S \setminus v}$, a degree-at-most-$2$ polynomial in
the other variables. Four pruning rules are used, each a necessary condition implied by
(E1)--(E3) on any genuine solution, so that discarding a branch violating one of them can never
lose a solution:
\begin{itemize}[leftmargin=1.4em]
  \item[(P1)] \emph{Budget}: the accumulated $\sum n_S^2$ never exceeds $16$, and $m_v \le 16$;
  this is (i).
  \item[(P2)] \emph{Bounded conditional values}: for a partial assignment of a subset $A$ of the
  link's variables, the terms confined to $A$ must keep the partial value of $L_v$ inside
  $[-4,4]$; this is (E1) evaluated at every completion.
  \item[(P3)] \emph{Weighted degree}: for a link variable $w$, $\sum_{T \ni w} |c_T| \le 4$;
  this is (E1) applied one derivative further, $D_w L_v \in \{0,\pm2,\pm4\}$, so
  $|a - b|, |a+b| \le 4$ forces $|a| + |b| \le 4$.
  \item[(P4)] \emph{Needs}: $\sum_{u \text{ unprocessed}} \max(0, 4 - m_u) \le 3\,(16 -
  \text{used})$, since each unprocessed vertex needs mass at least $4$ by (E2) and any future
  coefficient contributes its square to at most $3$ vertices.
\end{itemize}
Symmetry is broken by a canonical-form test using only (E3): the global sign is fixed once at
the top; fresh variables are introduced in increasing label order; and among all images of the
current partial solution under the group generated by permutations of the fresh used variables,
negations of fresh variables, and (when the fixed part of the link vanishes) global negation of
the link, only the lexicographically minimal image is kept. Every genuine solution has some
image satisfying all of these simultaneously, so rejecting non-canonical images discards only
redundant representatives of an orbit. Consequently the raw solution counts below are already
\emph{orbit} counts. Condition (ii) is never pruned: at the last vertex the program evaluates
$4 f_N(x)$ at all $2^n$ points and requires the value to lie in $\{+4,-4\}$; this full, unpruned
check is the only place a branch is accepted.

The $n = 11$ run was split into six shards by the residue of the index of the canonical
top-level link modulo $6$. Every shard enumerates the same full list of $368$ canonical
top-level links and recurses only on its own residue class, so the six shards partition the
search exhaustively and disjointly; the enumeration order is deterministic because the program
is single-threaded and depends only on $n$.

\begin{table}[ht]
\caption{Exhaustive canonical search for $F(n)$. ``Solutions'' are orbit counts under the
symmetry reduction described above.}\label{tab:search}
\centering
\begin{tabular}{rrrrrl}
\toprule
$n$ & solutions (orbits) & leaves & nodes & canonical rejects & time \\
\midrule
$3$  & $14$   & $26$     & $1{,}887$      & $176$        & $0.00$\,s \\
$4$  & $102$  & $225$    & $24{,}373$     & $2{,}755$    & $0.01$\,s \\
$5$  & $809$  & $2{,}329$ & $280{,}282$   & $31{,}155$   & $0.15$\,s \\
$6$  & $1{,}773$ & $6{,}538$ & $2{,}986{,}948$ & $187{,}681$ & $2.44$\,s \\
$7$  & $1{,}060$ & $11{,}786$ & $23{,}632{,}808$ & $530{,}408$ & $33.20$\,s \\
$8$  & $213$  & $4{,}015$ & $77{,}313{,}447$ & $1{,}288{,}595$ & $171.77$\,s \\
$10$ & $\mathbf{1}$ & $120$ & $44{,}447{,}494$ & $1{,}063{,}032$ & $824.60$\,s \\
$11$ (6 shards) & $\mathbf{0}$ & $16$ & $96{,}822{,}409$ & $3{,}747{,}449$ & $\approx 2.1$\,CPU-h \\
\bottomrule
\end{tabular}
\end{table}

Table~\ref{tab:search} records the results. The case $n = 9$ was not run: it is not needed for
the $F(10)$ and $F(11)$ claims. The value $n = 10$ giving exactly one orbit matches the known
truth, and the recovered solution is $16$ terms of coefficient $\pm1$, i.e.\ exactly $\XiC_3$ up
to the symmetry group; this is a strong sanity check on the canonical-form machinery, which is
the most intricate part of the program. The value $n = 11$ giving $0$ across all six disjoint
shards is the search's verdict that $R_3 \le 10$.

Every emitted solution at $n = 6, 7, 8, 10$ was re-checked by a from-scratch independent
verifier which parses the solution, evaluates $f_N$ by brute force at all $2^n$ points, and
tests (i), (ii) and genuine relevance (two points differing only in coordinate $i$ with
different values of $f_N$, rather than mere appearance in the support): $3047$ solutions
checked, $0$ failures. So no false positives are reported. We note the residual caveat
explicitly: the soundness argument above shows the \emph{pruning rules} cannot lose a solution,
but it is not a line-by-line proof that the C implementation is correct, and a false negative at
$n = 11$ caused by an implementation bug cannot be excluded by rerunning the same binary. An
earlier version of the program had a genuine bug of exactly this kind (scratch buffers indexed
by recursion depth only, hence aliased across vertex steps); it reported $0$ solutions and $2$
leaves at $n = 10$, where the corrected program reports $1$ solution and $120$ leaves. All
figures reported here are from the corrected program, and the $n = 10$ agreement with $\XiC_3$
is the check that detects that failure mode.

\subsection{Independent structure search}
Independently of both the hand proof and the search above, a referee ran a CDCL SAT search
(CaDiCaL via \texttt{pysat}) with blocking clauses over the \emph{term structures} of an
eleven-variable solution: the supports satisfying the link-consistency conditions, enumerated up
to isomorphism, with the Boolean sign vectors on each enumerated afterwards. The complete run
found exactly three isomorphism classes of link-consistent $16$-term supports on $11$ vertices
and \emph{zero} Boolean sign vectors on them:
\begin{itemize}[leftmargin=1.4em]
  \item degree sequence $4^{11}$, four quadratic and twelve cubic terms --- the $\delta = 4$
  case with weight partition $\{2,2\}$ and an (L4) link, plus an octahedron;
  \item degree sequence $4^{11}$, two linear and fourteen cubic terms --- the $\delta = 4$ case
  with weight partition $\{1,1,1,1\}$ and $|N| = 4$, plus an octahedron;
  \item degree sequence $(8, 4^{10})$, sixteen cubic terms --- the $\delta = 0$ case with link
  $C_4 \cup C_4$, i.e.\ two octahedra glued at a vertex.
\end{itemize}
These are precisely the three branches that the hand proof closes with the disjoint-sum lemma
(Lemma~\ref{lem:C}); every other branch is closed at the link level in
Appendix~\ref{app:cases} and is reported infeasible by the search. A later rerun restricted to
the branch where the distinguished vertex has mass $4$ reports $420 + 585$ labelled structures,
$2$ isomorphism classes and $0$ Boolean sign vectors, consistent with the above; that rerun's
log ends inside the mass-$6$ branch without a totals line, so we quote the complete run for the
totals.

\subsection{Partial formalisation in Lean 4}\label{sec:lean}
The finite statement $F(n)$ was transcribed into Lean~4 (toolchain
\texttt{leanprover/lean4:v4.23.0}) as a frozen definition, and the first structural steps toward
$F(12)$ were proved against it. The following is the exact status.

\emph{What is machine-checked.} Twenty-three theorems are proved and kernel-checked, with no
\texttt{sorry}, no \texttt{admit}, no \texttt{native\_decide}, no user axioms and no
\texttt{unsafe}/\texttt{implemented\_by}/\texttt{extern} escape hatches; an automated gate runs
\texttt{lake build}, greps for all of those, and runs \texttt{\#print axioms} on every listed
theorem. Every one of the twenty-three depends only on the three standard Lean axioms
\texttt{propext}, \texttt{Classical.choice}, \texttt{Quot.sound}. They cover:
\begin{itemize}[leftmargin=1.4em]
  \item the basic identities ($f_N^2 = 16$ in integer form, the pointwise membership (E1),
  divisibility and lower bounds on the mass, $\sum_i m_i \le 48$);
  \item Lemma~\ref{lem:masses} for $n \ge 11$ (mass in $\{4,6,8\}$), and that at least nine of
  the eleven vertices have mass $4$;
  \item the bookkeeping identity \eqref{eq:bookkeeping} in the form $e + \delta = 48 - 4n$;
  \item the first half of Lemma~\ref{lem:links}(a): four distinct link sets of size at most $2$,
  coefficients $\pm1$, symmetric difference empty, sign product $+1$, and every vertex in an
  even number of the four sets;
  \item Step 1 of Theorem~\ref{thm:no12} for $n = 12$ (all masses $4$, all terms cubic, all
  coefficients $\pm1$, all link sets of size $2$);
  \item Step 2 of Theorem~\ref{thm:no12}: four distinct pairs with empty symmetric difference
  form a $4$-cycle, hence for $n = 12$ every vertex link is a $4$-cycle.
\end{itemize}
The strongest certified statement is: in any $12$-variable solution of the frozen $F$, at every
vertex $v$ there are four distinct vertices $p,q,r,s$ such that every support set at $v$, with
$v$ removed, is one of the pairs $pq, qr, rs, sp$.

\emph{What is not machine-checked.} Neither $F(12)$ nor $F(11)$ is formalised. The remaining
steps for $F(12)$ are the octahedral closure (Step 3 of Theorem~\ref{thm:no12}) and the
disjoint-sum contradiction (Step 4); for $F(11)$, the whole of Appendix~\ref{app:cases}. Mathlib
has no classification of closed surfaces, so the closure arguments would have to be replaced by
their topology-free finite versions, and the exhaustive search of the previous subsection is far
beyond kernel \texttt{decide}. Theorem~\ref{thm:main-ns} is not formalised at all.

\section{Concluding remarks and open problems}\label{sec:open}

\subsection{The state of the small-degree values}
We now know
\[
  R_1 = 1, \qquad R_2 = 4, \qquad R_3 = 10, \qquad 22 \le R_4 \le 31,
\]
where $R_4 \ge 22 = 3 \cdot 2^3 - 2$ comes from $\XiC_4$ and $R_4 \le 31 = 4 \cdot 2^3 - 1$ from
Theorem~\ref{thm:main-ns}. Determining $R_4$ is the obvious next question. The gap is wide, and
the hypergraph in question has $64$ terms rather than $16$, so a direct repetition of
Appendix~\ref{app:cases} is not realistic; a new structural input is needed.

\begin{question}
What is $R_4$? Is it $22$?
\end{question}

\subsection{The extremality conjecture}
Since $R_2 = 4 = 3 \cdot 2^1 - 2$ and $R_3 = 10 = 3 \cdot 2^2 - 2$, both equal to the size of
the Chiarelli--Hatami--Saks construction, it is natural to conjecture that the construction is
always extremal.

\begin{conjecture}\label{conj:chs}
$R_d = 3 \cdot 2^{d-1} - 2$ for every $d \ge 2$; equivalently, $C_{44} = 3/2$.
\end{conjecture}

Conjecture~\ref{conj:chs} would close entry~44a of \cite{TaoOP44a} completely. We stress that
\cite{CHS20} supply the construction but do not state this conjecture; the evidence for it is
the two exact values above, and it is far from the known upper bound $4.394 \cdot 2^d$.

One natural line of attack is the recursion behind the construction: $\XiC_d$ is built from two
copies of $\XiC_{d-1}$ and two new selector variables, giving
$R_d \ge 2 R_{d-1} + 2$. A matching upper bound
\begin{equation}\label{eq:recursion}
  R_d \;\le\; 2 R_{d-1} + 2
\end{equation}
would prove Conjecture~\ref{conj:chs} by induction from $R_2 = 4$. The link structure of
Section~\ref{sec:ns} is the natural tool: Remark~\ref{rem:A} says that the derivative at any
variable of minimal influence $2^{1-d}$ is a character times the indicator of an affine subspace
of codimension $d-1$, \emph{without} any homogeneity hypothesis. What is missing is the
corresponding structure for variables of the next influence levels $2^{2-d}, 2^{3-d}, \dots$ ---
in $\XiC_d$ these are exactly the selector variables, and they are the variables that
\eqref{eq:recursion} would have to control.

\begin{question}
Describe the derivatives $D_i f$ of a degree-$d$ Boolean function at variables of influence
$2^{2-d}$. Does a structure theorem at the first few influence levels imply
\eqref{eq:recursion}?
\end{question}

\subsection{Uniqueness of the extremal degree-3 function}
The exhaustive search of Section~\ref{sec:comput} found exactly one orbit at $n = 10$, namely
$\XiC_3$, under the group generated by permutations of variables, negation of variables and the
global sign. So, computationally, the extremal degree-$3$ function is unique up to those
symmetries. We have no proof of this by hand.

\begin{question}
Prove without computer search that $\XiC_3$ is, up to permutation and negation of variables and
global sign, the unique degree-$3$ Boolean function with $10$ relevant variables. More
generally, is the extremal function unique at every $d$?
\end{question}

\subsection{On the method}
The reusable content of this paper is Lemma~\ref{lem:A} together with Remark~\ref{rem:A}: a
$\bits$-valued function whose Fourier mass is as small as granularity permits is forced to be a
character times an affine-subspace indicator. This is a statement about sparse granular
functions, with no degree hypothesis, and it is what converts an influence-counting argument
into a rigid combinatorial one. At $d = 3$ the resulting objects are closed triangulated
surfaces, which is why a complete classification is possible; for general $d$ the same device
yields cross-polytopes and a clean disjointness contradiction, which is exactly enough to strip
one variable off the Nisan--Szegedy bound, but not obviously more.

\section*{Acknowledgements}
The mathematics was developed with AI assistance (Claude). Every claim in
Sections~\ref{sec:ns}--\ref{sec:partII} and Appendix~\ref{app:cases} was reviewed by an
independent referee instance which supplied the corrections now incorporated, in particular the
fourth four-set link shape (L4) of Lemma~\ref{lem:links}(a) and the Euler-characteristic
Fact~\ref{fact:T}.

\appendix

\section{The eleven-variable case analysis}\label{app:cases}

This appendix proves Theorem~\ref{thm:no11}. Notation is that of Section~\ref{sec:partII}:
$f \colon \cube^{11} \to \cube$ has degree at most $3$ with all $11$ variables relevant,
$n_S = 4\Fhat{f}(S) \in \bbZ$, $\sum_S n_S^2 = 16$, $m_i = \sum_{S \ni i} n_S^2$, and by
\eqref{eq:bookkeeping} we have $e + \delta = 4$ with $\delta \in \{4,2,0\}$. We write ``$m = 4$
vertex'' for a vertex of mass $4$, and use the link shapes (L1)--(L4) of
Lemma~\ref{lem:links}(a) throughout. We refer to a term $S$ with $|S| = 3$ as a \emph{cubic},
$|S| = 2$ as a \emph{quadratic}, $|S| = 1$ as a \emph{linear} term.

\subsection{The case \texorpdfstring{$\delta = 4$}{delta = 4}: all masses are 4}
\label{app:delta4}

Here every $m_i = 4$, so every link is (L1), (L2), (L3) or (L4). The lower-order weight $4$ is
partitioned as one of $\{1,1,1,1\}$, $\{4\}$, $\{2,1,1\}$, $\{2,2\}$, $\{3,1\}$, with
contributions as listed in Section~\ref{sec:bookkeep}.

\subsubsection*{Partition $\{4\}$} A quadratic term with $|n_S| = 2$. By
Lemma~\ref{lem:links}(d) its endpoints would need mass $8$, contradicting $m_i = 4$.
\hfill$\times$

\subsubsection*{Partition $\{2,2\}$} Two linear terms $\{i\}$ and $\{j\}$, and no quadratics.
The link of $i$ contains $\emptyset$, and there are no quadratics, so it is (L4): the cubics at
$i$ are $\{i,a,b\}, \{i,b,c\}, \{i,a,c\}$. Likewise $j$ has an (L4) link.

If $j \in \{a,b,c\}$, say $j = a$, then the link of $a$ consists of $\emptyset$ together with a
triangle containing the pairs $\{i,b\}$ and $\{i,c\}$, which forces the cubic $\{a,b,c\}$. But
then the link of $b$ contains $\{i,a\}, \{i,c\}, \{a,c\}$, a triangle, although $b$ has no
lower-order term and $m_b = 4$, so the link of $b$ must be a $C_4$ --- and a $C_4$ contains no
triangle. \hfill$\times$

Hence $j \notin \{a,b,c\}$, and $a, b, c$ have (L1) links. The link of $a$ contains
$\{i,b\}$ and $\{i,c\}$, which completes to the $4$-cycle $i$--$b$--$w$--$c$, giving the cubics
$\{a,b,w\}$ and $\{a,c,w\}$; the same argument at $b$ gives $\{a,b,w'\}, \{b,c,w'\}$ and at $c$
gives $\{a,c,w''\}, \{b,c,w''\}$, and the shared cubics force $w = w' = w''$. Then the link of
$w$ contains the triangle $\{a,b\},\{a,c\},\{b,c\}$, so $w$ is (L4), so $w = j$ (the only other
vertex with a linear term). The terms
\[
  \{i\},\ \{j\},\ \{i,a,b\},\ \{i,b,c\},\ \{i,a,c\},\ \{j,a,b\},\ \{j,b,c\},\ \{j,a,c\}
\]
exhaust the vertices $i, j, a, b, c$ (each of mass $4$) and carry weight $8$; the remaining
weight $8$ lives on the other $6$ variables. So $f$ is a sum of two non-constant functions on
disjoint variable sets, impossible by Lemma~\ref{lem:C}. \hfill$\times$

\subsubsection*{Partition $\{3,1\}$} A constant term and one quadratic $\{i,a\}$. The link of
$i$ contains the singleton $\{a\}$, so it is (L2) or (L3) by
Lemma~\ref{lem:links}(a); but (L2) needs a linear term at $i$, and there is none (the weight is
spent on the constant and the quadratic), while (L3) needs two quadratics at $i$, and there is
only one quadratic in total. \hfill$\times$

\subsubsection*{Partition $\{2,1,1\}$} One linear term $\{i\}$ and two quadratics. The link of
$i$ contains $\emptyset$, so it is (L2) or (L4).

\emph{If it is} (L2) $= \{\emptyset, \{a\},\{b\},\{a,b\}\}$: the two quadratics are $\{i,a\}$
and $\{i,b\}$, and $\{i,a,b\}$ is a cubic. The link of $a$ contains the singleton $\{i\}$ (from
the quadratic $\{i,a\}$), so $a$ is (L2) or (L3). It cannot be (L2), which would need a linear
term $\{a\}$, while the only linear term is $\{i\}$. So it is (L3), of the form
$\{\{p\},\{q\},\{p,r\},\{q,r\}\}$ with $\{i\} = \{p\}$; the pair $\{i,b\}$ (coming from the cubic
$\{i,a,b\}$) must be $\{p,r\}$, so $r = b$, and then the fourth member $\{q,r\}$ requires a
quadratic $\{a,q\}$ with $q \ne i$. But the only quadratics are $\{i,a\}$ and $\{i,b\}$.
\hfill$\times$

\emph{If it is} (L4): no quadratic touches $i$. So every endpoint $u$ of a quadratic has
$u \ne i$, has a singleton in its link, and has no linear term of its own; hence $u$ is (L3) and
therefore lies on exactly two quadratics. The two quadratics would then form a $2$-regular
graph with $2$ edges, which does not exist. \hfill$\times$

\subsubsection*{Partition $\{1,1,1,1\}$} Four unit quadratics forming a graph $Q$, with no
linear and no constant term. Every endpoint $u$ of a quadratic has a singleton in its link,
hence is (L3) (it cannot be (L2), which needs a linear term), hence lies on exactly two
quadratics. So $Q$ is a $2$-regular graph with $4$ edges, i.e.\ a $4$-cycle
$i$--$a$--$j$--$b$--$i$. By (L3) at each of $i,a,j,b$, the two cubics at $i$ are
$\{i,a,c_i\}, \{i,b,c_i\}$ for some $c_i$, and similarly at $a$, $j$, $b$.

Let $S_0 = \{i,a,j,b\}$ and let $N$ be the set of cubics meeting $S_0$. Each member of $N$ meets
$S_0$ in $2$ or $3$ vertices, since it contains one of $i,a,j,b$ together with a $Q$-neighbour of
that vertex. Each of $i,a,j,b$ lies in exactly two cubics, so
$\sum_{T \in N} |T \cap S_0| = 8$. The seven vertices outside $S_0$ all have mass $4$, hence
\[
  28 = \sum_{T \in N} |T \setminus S_0| + 3M ,
\]
where $M$ is the number of cubics disjoint from $S_0$, while counting all $16$ terms gives
$16 = 4 + |N| + M$. Writing $k := \#\{T \in N : |T \setminus S_0| = 1\}$ (the others have
$|T \setminus S_0| = 0$), we get $3(12 - |N|) = 28 - k$, so
\[
  k = 3|N| - 8 \le |N| ,
\]
which forces $|N| \in \{3,4\}$.

\smallskip\noindent\emph{Sub-case $|N| = 4$.} All four cubics of $N$ have exactly two vertices in
$S_0$ and one outside. Take $\{i,a,c\} \in N$. The vertex $c$ has an (L1) link containing the
pair $\{i,a\}$, so the $4$-cycle in $\lk(c)$ is $i$--$a$--$y$--$z$--$i$, giving cubics
$\{c,a,y\}, \{c,y,z\}, \{c,z,i\}$. Now $a$'s only other cubic is $\{a,j,c_a\}$, so $y = j$ and
$c_a = c$; and $i$'s only other cubic is $\{i,b,c_i\}$, so $z = b$ and $c_i = c$; and
$\{c,j,b\}$ is a cubic. Thus
\[
  N = \big\{ \{c,a,i\},\ \{c,a,j\},\ \{c,j,b\},\ \{c,b,i\} \big\},
\]
$m_c = 4$ is exhausted, and $i, a, j, b$ are exhausted too (two quadratics and two cubics each).
Hence no other term meets $\{i,a,j,b,c\}$, and
$f = g(x_i,x_a,x_j,x_b,x_c) + h(\text{the other } 6 \text{ variables})$ with both parts
non-constant, impossible by Lemma~\ref{lem:C}. \hfill$\times$

\smallskip\noindent\emph{Sub-case $|N| = 3$.} Then $k = 1$: two cubics lie inside $S_0$ and one
meets it in an edge. Two $3$-subsets of the $4$-cycle's vertex set share two vertices, so up to
the symmetry of the $4$-cycle they are $\{i,a,j\}, \{i,a,b\}$ (sharing the $Q$-edge $\{i,a\}$) or
$\{i,a,j\}, \{j,b,i\}$ (sharing the diagonal $\{i,j\}$).
\begin{itemize}[leftmargin=1.4em]
  \item First: $\lk(i) = \{\{a\},\{b\},\{a,j\},\{a,b\}\}$, which is not of the form
  $\{\{p\},\{q\},\{p,r\},\{q,r\}\}$ required by (L3). \hfill$\times$
  \item Second: $\lk(a)$ contains the pair $\{i,j\}$ (from the cubic $\{i,j,a\}$); but an (L3)
  link at $a$ has the form $\{\{p\},\{q\},\{p,r\},\{q,r\}\}$ with $\{p,q\} = \{i,j\}$ and
  $r \notin \{i,j\}$, so $\{i,j\}$ would have to be $\{i,r\}$, forcing $r = j$.
  \hfill$\times$
\end{itemize}

So $\delta = 4$ is impossible.

\subsection{The case \texorpdfstring{$\delta = 2$}{delta = 2}: one vertex of mass 6}
\label{app:delta2}

Here exactly one vertex $v$ has $m_v = 6$ and the other ten have mass $4$. The lower-order
weight $2$ is either one linear term $\pm1$ or two unit quadratics.

\subsubsection*{One linear term $\{l\}$}
The link of $l$ contains $\emptyset$. There are no quadratics, so every vertex other than $l$
has a link made of pairs only: a $C_4$ if its mass is $4$, and (for mass $6$) a simple graph
with $6$ edges and all degrees even, i.e.\ $C_6$, two disjoint triangles, or a bowtie (two
triangles sharing a vertex).

\emph{Suppose $m_l = 4$.} Then $\lk(l)$ is (L4), giving cubics
$\{l,a,b\}, \{l,b,c\}, \{l,a,c\}$.

Assume first $v \notin \{a,b,c\}$, so $a,b,c$ have $C_4$ links. The link of $a$ contains
$\{l,b\}$ and $\{l,c\}$, which completes to $l$--$b$--$w$--$c$, giving cubics
$\{a,b,w\}, \{a,c,w\}$; similarly at $b$ and at $c$, and the shared cubics force a common $w$,
so the cubics $\{a,b,w\}, \{a,c,w\}, \{b,c,w\}$ all exist. Then $\lk(w)$ contains the triangle
$abc$, so it is not a $C_4$; hence $w = v$ (and $w \ne l$, since $l$'s link is already
exhausted). So $\lk(v)$ contains the triangle $abc$ and $m_v = 6$: thus $\lk(v)$ is two
triangles $abc \cup xyz$, or a bowtie. A bowtie containing $abc$ has its centre in
$\{a,b,c\}$, say $a$, giving $\lambda_{va} = 4$ and hence a vertex of degree $4$ in $\lk(a)$,
which is not a $C_4$ --- contradiction. So $\lk(v) = abc \cup xyz$, with cubics
$\{v,x,y\}, \{v,y,z\}, \{v,x,z\}$. Now $x, y, z$ have $C_4$ links, and the same completion
argument yields cubics $\{x,y,t\}, \{x,z,t\}, \{y,z,t\}$ for a common $t$, so $\lk(t)$ contains
a triangle; but $t \notin \{v, l\}$ (both of those links are already exhausted), so $t$ has a
$C_4$ link containing a triangle. \hfill$\times$

Assume now $v \in \{a,b,c\}$, say $v = a$. Then $b$ and $c$ have $C_4$ links. The link of $b$
contains $\{l,a\}$ and $\{l,c\}$, giving cubics $\{a,b,w\}, \{b,c,w\}$; the link of $c$ contains
$\{l,a\}$ and $\{l,b\}$, giving $\{a,c,w'\}, \{b,c,w'\}$; the shared cubic $\{b,c,\cdot\}$ forces
$w = w'$. So $\lk(w)$ contains the triangle $abc$. Here $w \notin \{l, a\}$ (taking $w = l$ would
repeat the term $\{l,a,b\}$, and $w = a$ is not allowed since $\{a,a,b\}$ is not a set), so $w$
has a $C_4$ link containing a triangle. \hfill$\times$

\emph{Hence $m_l = 6$, i.e.\ $l = v$.} Then $\lk(v)$ is $\emptyset$ together with five pairs
having all degrees even, i.e.\ a $5$-cycle $a_1 a_2 a_3 a_4 a_5$, with cubics
$\{v, a_k, a_{k+1}\}$ (indices mod $5$). Each $a_k$ has an (L1) link containing
$\{v, a_{k-1}\}$ and $\{v, a_{k+1}\}$, which completes to a $4$-cycle
$v$--$a_{k+1}$--$w_k$--$a_{k-1}$--$v$, giving cubics $\{a_k, a_{k\pm1}, w_k\}$. If $w_k$ were
another $a_p$, some $a$'s link would contain a triangle (for instance $w_1 = a_3$ makes $\lk(a_2)$
contain $\{v,a_1\}, \{v,a_3\}, \{a_1,a_3\}$), impossible inside a $C_4$. So $w_k$ lies outside
$\{v, a_1, \dots, a_5\}$. The cubic $\{a_k, a_{k+1}, w_k\}$ is one of $a_{k+1}$'s two non-$v$
cubics, so $w_{k+1} = w_k$; hence all $w_k$ are equal to a single vertex $w$, which then lies in
five cubics, i.e.\ $m_w = 5$, contradicting Lemma~\ref{lem:masses}. \hfill$\times$

\subsubsection*{Two quadratics}
An endpoint $u$ of a quadratic with $m_u = 4$ has a singleton in its link, so its link is (L2) or
(L3); (L2) needs a linear term, and there is none, so it is (L3), which needs two quadratics at
$u$: say $\{u,p\}$ and $\{u,q\}$ with $p \ne q$. If $p \ne v$ then $p$ is a mass-$4$ endpoint too
and would need two quadratics at $p$, but both quadratics already meet $u$ and there are only
two in total --- so $p = v$, and likewise $q = v$, contradicting $p \ne q$. So no mass-$4$ vertex
is an endpoint of a quadratic, which means each quadratic has both endpoints equal to $v$, absurd.
\hfill$\times$

So $\delta = 2$ is impossible.

\subsection{The case \texorpdfstring{$\delta = 0$}{delta = 0}: a homogeneous cubic}
\label{app:delta0}

Now $f = \frac14\sum_S n_S \chi_S$ is supported on triples, with $\sum_S n_S^2 = 16$, and the
degree sequence $(m_i)$ is either $(8, 4^{10})$ or $(6,6,4^9)$. No $|n_S| = 2$ occurs: by
Lemma~\ref{lem:links}(d) all three vertices of such a term would need mass $8$, and there is at
most one vertex of mass $8$. So $\mathcal{F}$ is a $3$-uniform hypergraph with $16$ triples.
Links are: a $C_4$ when $m = 4$ (Lemma~\ref{lem:links}(a)); for $m = 6$, a simple graph with $6$
edges and all degrees even, i.e.\ $C_6$, two disjoint triangles, or a bowtie; for $m = 8$, a
simple graph with $8$ edges and all degrees even.

\subsubsection{Sub-case $(8, 4^{10})$}
Let $v$ be the vertex with $m_v = 8$. If some $x$ had $\lambda_{vx} = 4$, then $v$ would have
degree $4$ in $\lk(x)$, so $\lk(x)$ would not be a $C_4$, contradicting $m_x = 4$ (and
$x \ne v$, since $x$ is a vertex of $\lk(v)$). Hence $\lk(v)$ is $2$-regular: $C_8$,
$C_3 \cup C_5$, or $C_4 \cup C_4$; and every pair lies in $0$ or $2$ triples, so $|\mathcal{F}|$
is a closed pseudo-manifold, with all links single cycles except possibly at $v$.

\paragraph{$\lk(v) = C_8$.} Then $|\mathcal{F}|$ is a closed surface with $V = 11$, $E = 24$,
$F = 16$, so $\chi = 3$. Consider its components. A closed surface with all degrees $4$ except
at most one vertex of degree $8$ has degree sum $4V' + 4\cdot\mathbf{1}\{v \in \text{it}\}$. A
sphere needs degree sum $6V' - 12$, giving either $V' = 6$ (the octahedron, not containing $v$)
or $6V' - 12 = 4V' + 4$, i.e.\ $V' = 8$ --- a sphere on $8$ vertices with degrees $(8, 4^7)$,
which is impossible since $v$ would need $8$ neighbours among $7$ other vertices. A projective
plane needs $6V' - 6$, giving $V' = 3$ or $V' = 5$, both below the minimum of $6$ vertices. And
$\chi' \le 0$ is excluded by Fact~\ref{fact:T}. So $\chi = 3$ cannot be realised as a sum of
component characteristics. \hfill$\times$

\paragraph{$\lk(v) = C_4 \cup C_4$.} Split $v$ into two vertices $v_1, v_2$, each of degree $4$.
The result is a closed surface with $12$ vertices all of degree $4$, so $\chi = 12 - 24 + 16 = 4$
and, exactly as in Step 3 of Theorem~\ref{thm:no12}, it is two octahedra. If $v_1$ and $v_2$ lay
in the same octahedron --- this sub-case is in fact empty, since the two $C_4$'s of $\lk(v)$ are
vertex-disjoint while two vertices of an octahedron share two or four neighbours --- the other
octahedron would use $6$ variables disjoint from everything else and $f$ would be a disjoint
sum, impossible by Lemma~\ref{lem:C}. Otherwise the two octahedra are glued at the identified
vertex $v$. Write
\[
  f = g_1(x_v, y) + g_2(x_v, z), \qquad y, z \text{ disjoint } 5\text{-tuples}.
\]
Fixing $x_v = 1$ gives $g_1(1,y) + g_2(1,z) \in \{\pm1\}$ for all $y, z$, so by
Lemma~\ref{lem:C} one of them is constant. But $g_1(1,y)$ is $\frac14$ times the sum of the eight
face-characters of the first octahedron with $x_v = 1$ substituted, which are eight distinct
non-constant monomials in $y$ (four quadratic and four cubic), so $g_1(1,\cdot)$ is not constant;
and the same holds for $g_2(1,\cdot)$. \hfill$\times$

\paragraph{$\lk(v) = C_3 \cup C_5$.} Split $v$ into $v_1$ of degree $3$ and $v_2$ of degree $5$:
a closed surface with $12$ vertices, degrees $(3,5,4^{10})$, and $\chi = 4$. Components with
$\chi' \le 1$ are excluded: $\chi' \le 0$ by Fact~\ref{fact:T}; and a projective plane needs
degree sum $6V' - 6$, which with all degrees at most $5$ forces $V' \le 6$ (and a triangulated
projective plane has at least $6$ vertices, by Step 3 of Theorem~\ref{thm:no12}), while the only
$6$-vertex projective plane has all six degrees equal to $5$, whereas we have a single degree-$5$
vertex. Hence both components are spheres, with degree sum $6V' - 12$ and degrees drawn from
$\{3,4,5\}$:
\begin{itemize}[leftmargin=1.4em]
  \item containing both the $3$ and the $5$: $8 + 4a = 6(a+2) - 12$ gives $a = 4$, a $6$-vertex
  sphere with degrees $(5,4,4,4,4,3)$. The degree-$5$ vertex is adjacent to all the others, its
  link is a $5$-cycle, and the remaining $3$ faces triangulate that pentagon with two diagonals;
  any two diagonals of a pentagon share a vertex or cross, so some vertex receives diagonal
  degree $2$ or the triangulation is invalid, and the required degree pattern
  $(4,4,4,4,3)$ on the rim, i.e.\ diagonal degrees $(1,1,1,1,0)$, cannot occur.
  \item containing the $3$ but not the $5$: $3 + 4a = 6(a+1) - 12$ gives $a = 4.5$, not an
  integer.
  \item containing the $5$ but not the $3$: $5 + 4a = 6(a+1) - 12$ gives $a = 5.5$, not an
  integer.
\end{itemize}
\hfill$\times$

\subsubsection{Sub-case $(6,6,4^9)$}
Let $v$ and $v'$ be the two vertices of mass $6$.

\paragraph{Some link is a bowtie.} Say $\lk(v)$ is a bowtie with centre $w$. Then
$\lambda_{vw} = 4$, so $v$ has degree $4$ in $\lk(w)$, which forces $m_w = 6$ and $\lk(w)$ a
bowtie too, so $w = v'$ and, symmetrically, $\lk(v')$ is a bowtie centred at $v$. The four common
triples are $\{v,v',a_k\}$, $k = 1,\dots,4$. The bowtie at $v$ adds two further triples
$\{v,a_1,a_2\}, \{v,a_3,a_4\}$ (for some pairing of the $a_k$), and the one at $v'$ adds
$\{v',a_i,a_j\}, \{v',a_k,a_l\}$ for some pairing.
If the two pairings coincide, then $\lk(a_1)$ contains $\{v,v'\}, \{v,a_2\}, \{v',a_2\}$, a
triangle, hence is not a $C_4$. Otherwise, say $v$ pairs $(a_1a_2)(a_3a_4)$ while $v'$ pairs
$(a_1a_3)(a_2a_4)$: then $\lk(a_1)$ contains $\{v,v'\}, \{v,a_2\}, \{v',a_3\}$, and being a $C_4$
it forces the fourth triple $\{a_1,a_2,a_3\}$; but then
$\lk(a_2) = \{\{v,v'\},\{v,a_1\},\{v',a_4\},\{a_1,a_3\}\}$ has $a_3$ and $a_4$ of degree $1$, so
it is not a $C_4$. \hfill$\times$

\paragraph{No link is a bowtie.} Then every pair lies in $0$ or $2$ triples and $|\mathcal{F}|$
is a closed pseudo-manifold, singular only at $v$ or $v'$ and only when the link there is two
triangles.

\emph{If both links are cycles}, $|\mathcal{F}|$ is a closed surface with
$\chi = 11 - 24 + 16 = 3$. Component degree sums are
$4V' + 2\cdot\mathbf{1}\{v\} + 2\cdot\mathbf{1}\{v'\}$.
Spheres: $V' = 6$ (containing neither $v$ nor $v'$, the octahedron); or, containing exactly one
of them, $4V' + 2 = 6V' - 12$, i.e.\ $V' = 7$ with degrees $(6,4^6)$ --- here the degree-$6$
vertex is adjacent to all $6$ others, its link is a hexagon, and the remaining $4$ faces
triangulate the hexagon with $3$ diagonals, so the six rim vertices have degrees $3 + d_i$ with
$\sum_i d_i = 6$; all degrees equal to $4$ needs $d_i = 1$ for every $i$, a perfect matching of
the hexagon by pairwise non-crossing diagonals, which does not exist (every hexagon
triangulation has an ear). Containing both $v$ and $v'$: $4V' + 4 = 6V' - 12$, i.e.\ $V' = 8$
with degree sum $36$, two degree-$6$ vertices and six degree-$4$ vertices --- this is not
immediately absurd, but then the remaining components would have to be octahedra on the
remaining $3$ vertices, which is impossible; and $V' = 11$ as a single component gives
$\chi = 3 \ne 2$. Projective planes: $6V' - 6 = 4V' + 2s$ with $s \in \{0,1,2\}$ gives
$V' = 3, 4, 5$, all below $6$. And $\chi' \le 0$ is excluded by Fact~\ref{fact:T}. So $\chi = 3$
is not realisable. \hfill$\times$

\emph{If $\lk(v)$ is two triangles}, split $v$ into $v_1, v_2$ of degree $3$ each. If $\lk(v')$
is a cycle we obtain $12$ vertices with degrees $(3,3,6,4^9)$ and $\chi = 4$; both components
are then spheres, since $\chi' \le 0$ is excluded by Fact~\ref{fact:T} and a projective plane
would need $6V' - 6 \le 6 + 4(V'-1)$, i.e.\ $V' \le 4 < 6$. The sphere containing $v'$ satisfies
$6 + 3a + 4b = 6(1 + a + b) - 12$, i.e.\ $3a + 2b = 12$ with $a \le 2$: the solution
$(a,b) = (2,3)$ gives a $6$-vertex sphere with $v'$ of degree $6$, impossible as there are only
$5$ other vertices; and $(a,b) = (0,6)$ gives the $(6,4^6)$ sphere excluded above. \hfill$\times$

If $\lk(v')$ is also two triangles, both split, giving degrees $(3,3,3,3,4^9)$ with $\chi = 5$:
odd, hence impossible as a sum of sphere characteristics; projective planes with all degrees at
most $4$ need $6V' - 6 \le 4V'$, i.e.\ $V' \le 3 < 6$; and $\chi' \le 0$ is excluded by
Fact~\ref{fact:T}. \hfill$\times$

\medskip
So $\delta = 0$ is impossible, and with Sections~\ref{app:delta4} and \ref{app:delta2} all cases
are exhausted. This proves Theorem~\ref{thm:no11}. \hfill$\square$

\begin{thebibliography}{99}

\bibitem{CHS20}
J.~Chiarelli, P.~Hatami and M.~Saks,
\emph{An asymptotically tight bound on the number of relevant variables in a bounded degree
Boolean function},
Combinatorica \textbf{40} (2020), 237--244.
\texttt{arXiv:1801.08564}.

\bibitem{NS94}
N.~Nisan and M.~Szegedy,
\emph{On the degree of Boolean functions as real polynomials},
Computational Complexity \textbf{4} (1994), 301--313.
\todo{page range not independently re-verified in this project.}

\bibitem{ODbook}
R.~O'Donnell,
\emph{Analysis of Boolean Functions},
Cambridge University Press, 2014.

\bibitem{TaoOP44a}
T.~Tao et al.,
\emph{Optimization problems}, entry 44a (relevant variables in bounded-degree Boolean
functions).
\url{https://teorth.github.io/optimizationproblems/constants/44a.html}.
Accessed 2026-09-09.

\bibitem{Wel19}
J.~Wellens,
\emph{Relationships between the number of inputs and other complexity measures of Boolean
functions},
preprint, \texttt{arXiv:1903.08214} (2019). Superseded by \cite{Wel22}.

\bibitem{Wel22}
J.~Wellens,
\emph{Relationships between the number of inputs and other complexity measures of Boolean
functions},
Discrete Analysis \textbf{2022}:19.
\texttt{arXiv:2005.00566}.

\end{thebibliography}

\end{document}
